\section{E. Classification and deformation questions}
\subsection{E.1. Higher logarithmic transformations}
We distinguish Viaclovsky's indices $0,1$ from the indices $1,2$ of
the question.  Put $m=3k$, $n=4l$, and assume $\gcd(m,n)=1$.
The lattice coordinates in this subsection are integral homology
coordinates, and matrices act on columns.

\begin{proposition}[E.1.1 (Proved): the integral dictionary]
Let $(\delta,a,b,d)$ be the marking in
\cite[Section 2, equations (2.1), (2.2), (2.7), (2.11)--(2.12)]{ViaLog26push},
and put $u_0=d+2a+2b$, $u_1=d+3a+3b$.
The homomorphism from the marking of the question to this marking is
\[
 P_V=\begin{pmatrix}
 -2&1&0&1\\-6&0&-1&0\\0&-1&0&0\\-1&0&0&0
 \end{pmatrix},\qquad \det P_V=-1.
\]
It intertwines the two finite monodromies and satisfies
\[
 P_V\widehat\delta=\delta,\qquad
 P_V\varepsilon=-u_0,\qquad
 P_V\varepsilon'=\delta-u_1.
\]
With the counterclockwise deck generators of that paper, its data
$m_0=3k$, $m_1=4l$ and
$a_i=[v_i/m_i]$ with $v_i=\rho_i\delta+N_i u_i$ correspond to
\[
 w_1=N_0\varepsilon-\rho_0\widehat\delta,
 \qquad
 w_2=N_1\varepsilon'-(\rho_1+N_1)\widehat\delta.
\]
For the section-preserving cusp comparison used there, $l_0=0$ and
$p=-(nN_0+mN_1)$.  Thus its sphere condition is
$b:=nN_0+mN_1=\pm1$.
The brackets in $a_i$ mean the torsion point in the smooth
root torus, using its section-preserving origin and the
transported period vector.  Replacing $v_i$ by an integral
multiple $m_i$ times a period does not change that torsion
point, although its chosen lift fixes the collar bookkeeping.
\end{proposition}
\begin{proof}
In Engel's marking the four columns of the change to Viaclovsky's
basis are
\[
 Q_E=\begin{pmatrix}-2&1&0&-2\\-1&0&0&0\\3&-1&1&0\\0&0&0&1\end{pmatrix}.
\]
Direct integral multiplication gives $P_V=Q_E^{-1}P_E$ and
$V_0P_V=P_VA_1$, $V_1P_V=P_VA_2$, where $V_0,V_1$ are the
matrices in that paper's Lemma 2.1.  Multiplication by the three
specified invariant vectors gives the displayed identities.
Inverting a clockwise affine generator replaces its invariant
translation $w_j/m_j$ by $-w_j/m_j$.  Consequently the actual
counterclockwise translation vector is $-P_Vw_j$, which gives the
two formulas, including the otherwise easily missed $N_1$ in the
second coefficient.  Viaclovsky's reference section extends through
the smooth locus of the cusp; its comparison therefore has zero
class in the quotient lattice of toric directions.  This is precisely
$l_0=0$.  Substitution into the definition of $p$ proves its formula.
These matrix identities are independently executed in
\texttt{code/e\_push\_exact.py}, with their output attached.
\end{proof}

For a general cusp collar $\sigma$, write $c(\sigma)$ for
its residual translation defect in
$\overline\Lambda=\Lambda/(\mathbb Z\widehat w+
\mathbb Z\widehat\delta)$, using the convention
$\sigma(g_0z)-R_0\sigma(z)=\Pi_c(g_0z)\ell(\sigma)$
and $c(\sigma)=[\ell(\sigma)]$.
Set $l_0=\gamma(c(\sigma))$.  In the question's notation
$w_j=r_j\varepsilon_j+s_j\widehat\delta$, the integer is
$p=mn l_0-nr_1-mr_2$.
The section-preserving source data of E.1.1 have $c(\sigma)=0$.

\begin{theorem}[E.1.2 (Proved): classification in Viaclovsky's coordinates]
Use a local logarithm $c$ of the transported linearization,
$u=e^{-2\pi i c}$, and the same reference section and cusp comparison
on both sides.  Two sphere members with data
$(m,\rho_0,N_0;n,\rho_1,N_1;c)$ and primed data are biholomorphic
if and only if their multiplicities agree and there are
$\alpha\in\{1,-1\}$ and $h=c'-c\in\mathbb Z$ such that
\[
 b'=\alpha b,\qquad
 hN'_0\equiv\rho'_0-\alpha\rho_0\pmod{k},\qquad
 hN'_1\equiv\rho'_1-\alpha\rho_1\pmod{2l}.
\]
Every such biholomorphism covers the identity of $B$.
For fixed discrete data the period identification subgroup is
$2kl\mathbb Z$.  At fixed $u$ all sphere translation data give
exactly $2kl$ classes.  A positive $u$-circuit acts on them by a
single cycle.
\end{theorem}
\begin{proof}
We first give the complete sufficiency argument in the question's
clockwise coordinates, writing $w_j=r_j\varepsilon_j+s_j\widehat\delta$.
Set $S_h=I+hE$ and $d_j=S_hw'_j-\alpha w_j$.
The repeated norm maps have images
\[
 N_m\Lambda=m\mathbb Z\varepsilon+k\mathbb Z\widehat\delta,
 \qquad
 N_n\Lambda=n\mathbb Z\varepsilon'+2l\mathbb Z\widehat\delta.
\]
This follows directly from
\[
 N_3v=3\gamma(v)\varepsilon+(2v_u+v_w+3v_\delta)\widehat\delta,
 \quad
 N_4v=4\gamma(v)\varepsilon'+(2v_u+2v_w+4v_\delta)\widehat\delta.
\]
Suppose $p'=\alpha p$ and
$hr'_1\equiv\alpha s_1-s'_1\pmod k$,
$hr'_2\equiv\alpha s_2-s'_2\pmod{2l}$.
Reduction of $p'=\alpha p$ modulo $m,n$ gives
$r'_j\equiv\alpha r_j\pmod{m_j}$.
We may choose integral solutions of $N_{m_j}\ell_j=-d_j$ explicitly:
\[
 \begin{split}
 \ell_1&=\left(\frac{\alpha r_1-r'_1}{m},0,
 -\frac{s'_1+hr'_1-\alpha s_1}{k},0\right)^t,\\
 \ell_2&=\left(\frac{\alpha r_2-r'_2}{n},
 -\frac{s'_2+hr'_2-\alpha s_2}{2l},0,0\right)^t.
 \end{split}
\]
Put $y_j=d_j/m_j+\ell_j$; these vectors belong to the rational
kernel of $N_{m_j}$ and have first coordinate zero.
The following rational vectors solve $(I-A_j)b_j=y_j$:
\[
 \begin{split}
 b_1&=(0,(2y_{1u}+y_{1w})/3,(y_{1w}-y_{1u})/3,0)^t,\\
 b_2&=(0,(y_{2u}-y_{2w})/2,(y_{2u}+y_{2w})/2,0)^t.
 \end{split}
\]
The last coordinates of the equations follow from the displayed
norm-kernel equations.  The maps
\[
 (s,\zeta)\longmapsto(s,\alpha\zeta+\Pi_c(s)b_j)
\]
are holomorphic maps of the normalized torus families over the
degree-$m_j$ root discs.  Their inverses are holomorphic, and period
covariance and $(I-A_j)b_j=d_j/m_j+\ell_j$ give exact commutation
with the prescribed affine deck generators, modulo the integral
period $\ell_j$.  They therefore descend to biholomorphisms of the
finite fillings, including the full multiple fibres.

On the punctured root collar the corresponding middle translation is
\[
 A_j=\Pi_cb_j+\frac{\log s_j}{2\pi i}\Pi_cd_j.
\]
Its integral defect is $\ell_j$.  The middle cocycle relation fixes
$\ell_0=-A_0(\ell_1+A_1\ell_2)$.  In the quotient
$\overline\Lambda=\Lambda/(\mathbb Z\widehat w+\mathbb Z\widehat\delta)$,
the first-coordinate condition is exactly $p'=\alpha p$.
Adding an integral multiple of $(0,1,-2,0)^t\in\ker N_m$ to
$\ell_1$ changes the other quotient coordinate by a generator,
without changing this first coordinate.  We can consequently arrange
\[
 [\ell_0]=-\alpha c(\sigma)+c(\sigma')\quad\hbox{in }\overline\Lambda.
\]
An integral toric vector then supplies $q\in\mathbb Z^2$ such that
the cusp reference $sq+\alpha\sigma-\sigma'$ has precisely defect
$\ell_0$.  Here $s=(2\pi i)^{-1}\log t_c$ is the cusp logarithm;
the maps of cusp toric covers are
\[
 (x,t_c)\longmapsto(t_c^q a(t_c)x^\alpha,t_c),
 \qquad a(0)\in(\mathbb C^*)^2.
\]
Integer translations of the cocharacter lattice, its inversion,
and multiplication by a holomorphic unit preserve the specified
hexagonal fan and all its strata.  These maps and their inverses
extend holomorphically across the three double curves and both
triple points.

We justify the global analytic correction.  Differences of allowed
finite extensions are the equivariant holomorphic vertical vectors
on the complete root discs.  The extra subgroup of order $k$ at
the first point acts trivially on the vertical vector coefficients
and rotates $s$; its invariant holomorphic coefficients are exactly
those in $s^k$.  The remaining order-three equivariance is the
original one.  The argument at the second point uses $s^l$ and
order four.  Thus the coarse difference sheaf is literally the same
sheaf $\mathcal V_B$ as for multiplicities $(3,4)$, not an assumed
analogue.  Its first-coordinate quotient is
$\mathcal O_B(-p_0)$: the modular frame has vanishing orders $2k,l$
on the two enlarged roots and a simple cusp pole.  Its fixed
second-coordinate subline is $\mathcal O_B$.  Finite averaging
provides local lifts, hence
\[
 0\longrightarrow\mathcal O_B\longrightarrow\mathcal V_B
 \longrightarrow\mathcal O_B(-p_0)\longrightarrow0.
\]
Both end terms have vanishing $H^1$ on $\mathbb P^1$, so
$H^1(B,\mathcal V_B)=0$.  The permitted local translations form a
torsor under this sheaf: on each overlap their defects agree, so
their difference is a single-valued section of the indicated
difference sheaf.  Vanishing of its \v Cech class gives a global
middle translation.  Its overlap equalities are exactly the two
finite formulas and the cusp formula just displayed.  Consequently
the local biholomorphisms glue, with their inverses, over $B$.

For necessity the base is intrinsic even when its anticanonical
pencil changes.  If $F:X\to X'$ is a biholomorphism, the pullback
of the Fubini--Study form by $f'F$ is exact because
$H^2(S^6,\mathbb R)=0$.  Its restriction to a regular complex torus,
wedged with a flat K\"ahler form, has zero integral.  Positivity
forces $f'F$ to be constant on that torus.  The proper analytic
image of $(f,f'F)$ is therefore a curve of degree one over $B$;
normalization and the inverse argument give a base automorphism.
Its three monodromy orders $3,4,\infty$ are distinct, so it fixes
all three marked points and is the identity.  Multiplicities agree.
On ordinary fibres the lift is affine.  The simultaneous rational
centralizer of $A_1,A_2$ is $\mathbb Q I+\mathbb Q E$, as follows
by solving their entrywise commuting equations.  Comparison of the
constant last two columns of $\Pi_c$ and unimodularity then gives
fibre-linear part $\alpha I$, integer difference $h$, and lattice
map $\alpha S_h^{-1}$.  The normalized root base change is
functorial, so the lifted map commutes with the actual deck
generator.  Summing its defect around that generator gives the
two norm conditions above.  A cusp map has an integral toric
winding: at a generic point of a toric component, its ratio of
characters is $t_c^q$ times a holomorphic unit.  Its class in
$\overline\Lambda$ vanishes.  The middle cocycle relation then
gives $p'=\alpha p$.  This proves all the necessary conditions.

Substitute E.1.1.  In the second congruence the extra term
$N'_1-\alpha N_1$ is divisible by $4l$, and hence drops out
modulo $2l$.  We obtain exactly the statement's two congruences.
For fixed data $b\ne0$ excludes $\alpha=-1$.
The invertibility of $N_0$ modulo $k$, of $N_1$ modulo $2l$, and
$\gcd(k,2l)=1$ give $h\in2kl\mathbb Z$.
For varying data the Chinese remainder theorem gives exactly one
shift class modulo $2kl$, and all residue classes of
$\rho_0\bmod k$, $\rho_1\bmod2l$ occur without changing $b$ or
freeness.  Finally a positive $u$-circuit sends $c$ to $c-1$.
Its translation on $\mathbb Z/(2kl)$ is a generator.
\end{proof}

The preceding theorem answers Question 3 of
\cite[Section 1.1]{ViaLog26push}.  The original return asserted the
same arithmetic classification in B.3.1; the contribution here is
the explicit dictionary and the complete enlarged-root and collar
argument.  We make no priority claim.

\begin{corollary}[E.1.3 (Proved): holomorphic dependence of the comparisons]
For fixed discrete data in E.1.2 and a fixed admissible shift $h$,
the constructed comparisons are holomorphic isomorphisms of families
over a sufficiently small parameter disc.  In particular, for sphere
data with multiplicities $(3,4)$ the identification with the standard
family uses the parameter
\[
 W=e^{-\pi i(c+s_2)}=(-1)^{s_2}w.
\]
The smooth monodromy conjugacy class around the positive parameter
circle is unchanged by this identification of twist data.
\end{corollary}
\begin{proof}
In E.1.2 the rational vectors $b_j$, integral defects $\ell_j$,
and toric exponents are independent of $c$.  All period matrices,
root comparisons, and cusp unit factors are holomorphic in $c$.
On a small Stein parameter disc $D$, the correction sheaf is
$\mathcal V_B\boxtimes\mathcal O_D$.  Its first cohomology on
$B\times D$ is zero: project to $D$, use the same two-chart
complex and $H^1(B,\mathcal V_B)=0$, and then Stein acyclicity.
Thus the affine correction can be chosen holomorphically in $c$.
For $(k,l)=(1,1)$ the only nonvacuous finite congruence is
$h\equiv s_2-s_2'\pmod2$, because the admissible $r_2'$ is odd.
Taking $s_2'=0$ gives the displayed parameter identification.
Two choices of these family comparisons differ fibrewise by an
automorphism of the standard member.  The full automorphism
calculation audited in E.3.2 identifies that group with $\mathbb C^*$;
it is connected.  On overlapping parameter arcs the transitions
are therefore isotopic to the identity.  Composing around a circle
changes the resulting diffeomorphism only by such an isotopy and
by conjugation at the chosen base fibre.  Multiplication of the
parameter by $(-1)^{s_2}$ preserves the positive orientation.
\end{proof}

\subsection{E.2. Tangent cohomology for arbitrary coprime multiplicities}
\begin{theorem}[E.2.1 (Proved): tangent cohomology and local miniversality]
For every sphere member with multiplicities $(3k,4l)$ and the
admissible logarithmic data of E.1, put $z=kl$ and $d=3k+4l$.
On $\mathfrak B_{k,l}=\mathbb P(3k,4l)$ the root lift
$\bar f:X\to\mathfrak B_{k,l}$ satisfies
\[
 \begin{split}
 \mathcal E&=\mathcal O\oplus\mathcal O(-z),\\
 R^q\bar f_*\mathcal O_X&=
 \mathcal O,\quad\mathcal O\oplus\mathcal O(-z),\quad\mathcal O(-z),\\
 K_X&=\bar f^*\mathcal O(z-d).
 \end{split}
\]
The tangent cohomology is $(h^0,h^1,h^2,h^3)(X,T_X)=(1,1,1,0)$.
With the groupoid derivative convention
$\kappa(\Phi)=(\partial_c\Phi)\Phi^{-1}$, the class of
$\partial/\partial c$ is $e\otimes[\gamma]$.  The class of
$\partial/\partial u$ is
$-(2\pi i u)^{-1}e\otimes[\gamma]$.  Both are nonzero.
The Kuranishi germ is a smooth reduced line, and the fixed-data
period family is locally miniversal.
\end{theorem}
\begin{proof}
The root object is defined by the actual divisor identities
$f^*(0)=mS_1$ and $f^*(1)=nS_2$, with $m=3k$, $n=4l$.
It gives a holomorphic map to the analytic root stack.  Its base
changes to the two root discs are the actual smooth torus
families; the quotient actions are free.  At the cusp its charts
are $q=z_0z_1z_2$, so the source is smooth and the map is flat,
of relative dimension two, with Gorenstein fibres.  Properness is
checked on the root atlas, where it is properness of the compact
torus families and the original compact cusp filling.

Use weighted coordinates $x,y$ of weights $m,n$, with coarse
coordinate $t=y^m/(y^m-x^n)$.
The two root divisors have bundles $\mathcal O(n)$ and
$\mathcal O(m)$, and the cusp has bundle $\mathcal O(mn)$.
The root-forgetting morphism
\[
 r:\mathfrak B_{k,l}\longrightarrow\mathfrak B_{1,1},
 \qquad (x,y)\longmapsto(x^l,y^k),
\]
uses the group homomorphism $\lambda\mapsto\lambda^{kl}$ on
the weighted quotient.  It is a holomorphic morphism of analytic
stacks over $B$, with local maps $s\mapsto s^k$ and
$s\mapsto s^l$.  In particular $r^*\mathcal O(1)=\mathcal O(z)$.
For completeness, the analytic Picard group of this root curve has
the presentation
$\langle H,R_1,R_2\mid mR_1=H,\ nR_2=H\rangle$:
after subtracting the two stabilizer characters, a line bundle
descends to the coarse projective line.  Its only relations are
the two defining root identities.  The Smith form and
$\gcd(m,n)=1$ identify this group with $\mathbb Z$, with
$H=mn\mathcal O(1)$, $R_1=n\mathcal O(1)$, and
$R_2=m\mathcal O(1)$.  This proves the analytic torsion-free
Picard assertion used below.
Translations have identity vertical derivative and act trivially
on structure-sheaf cohomology.  Thus the vertical bundle computed
from the linear period factors is $r^*(\mathcal O\oplus
\mathcal O(-1))$.  The original determinant calculation can also
be repeated: its quotient line $L$ has
$L^{12}=\mathcal O(-p_0)=\mathcal O(-12z)$, so
$L=\mathcal O(-z)$ because the Picard group is torsion-free.
The fixed vector gives a subline $\mathcal O$; its extension
splits because $H^1(\mathcal O(z))=0$.

The cusp normalization resolution has terms
$\mathcal O_{\mathrm{dP}_6}$, three copies of
$\mathcal O_{\mathbb P^1}$, and two point sheaves.
Its global differential is $(a,b,c)\mapsto
(a-b+c,a-b+c)$, giving cohomology $(1,2,1)$; its alternating cup
pairing in degree one is nondegenerate.  The degree-one cocycle
$\gamma$ extends across the cusp because it vanishes on the toric
period lattice.  Its edge values $(0,-1,-1)$ have nonzero cup
product with $(1,0,-1)$.  On finite root charts its class is the
nonzero constant lattice cocycle $\gamma$, invariant under every
deck transformation.  These are explicit comparisons of the
same class, so it gives a nowhere-zero subline in
$G=R^1\bar f_*\mathcal O_X$.
The base-change and relative-duality argument E.3.1 applies on
these charts.  It gives $R^2\bar f_*\mathcal O_X=L$ and
$\det G=L$.  The extension of $L$ by the specified
$[\gamma]$-subline splits by the same $H^1(\mathcal O(z))=0$.
This proves all direct-image formulas.

The logarithmic tangent sequence is
\[
 0\longrightarrow\mathcal V=\bar f^*\mathcal E
 \longrightarrow\mathcal T=T_X(-\log W)
 \longrightarrow\mathcal H=\bar f^*\mathcal O(d-12z)
 \longrightarrow0.
\]
Its last map is the relative differential on smooth root charts
and sends $\sum a_i z_i\partial_{z_i}$ to
$(\sum a_i)q\partial_q$ at the cusp.  It is surjective there.
Since $K_{\mathfrak B}=\mathcal O(-d)$, determinants give
$K_X=\bar f^*\mathcal O(z-d)$.

We give the weighted cohomology calculation explicitly; simply
replacing $1$ by $kl$ in an acyclicity assertion would be incorrect.
The two quotient charts $x\ne0$, $y\ne0$ are finite quotients
of affine lines, and their overlap is a finite quotient of a
punctured line.  Finite averaging and Stein acyclicity compute
coherent cohomology by this two-chart complex.  Its homogeneous
Laurent monomials give
\[
 \begin{split}
 h^0\mathcal O(a)&=\#\{(i,j)\in\mathbb Z_{\ge0}^2:mi+nj=a\},\\
 h^1\mathcal O(a)&=h^0\mathcal O(-a-d).
 \end{split}
\]
The second equality identifies the remaining monomials
$x^{-i-1}y^{-j-1}$; it does not require a positivity theorem.
The bundles $\mathcal O(-z),\mathcal O(-2z)$ are acyclic.
Indeed a possible degree-one monomial would give
$m(i+1)+n(j+1)=az$, with $a=1$ or $2$.
Coprimality forces $i+1=lA$, $j+1=kB$, hence
$3A+4B=a$ with $A,B\ge1$, which is impossible.
Positive degree has no degree-one monomials, so
$H^1\mathcal O(z)=0$ as used above.

Put $b=d-12z$.  Both $b$ and $b-z$ are negative.
Degree-one monomials for $\mathcal O(b)$ would satisfy
$3A+4B=12$, with $A,B\ge1$, and none exist.
For $\mathcal O(b-z)$ the equation is $3A+4B=13$;
its only positive solution is $(A,B)=(3,1)$, producing
$x^{3l-2}y^{k-2}$ in the dual description.  It is allowed
exactly when $k>1$; write $e_k=1$ in that case and $e_k=0$
otherwise.  Projection formula and the two-column Leray sequence
now give
\[
 h^\bullet(\mathcal V)=(1,1,0,0),\qquad
 h^\bullet(\mathcal H)=(0,0,e_k,e_k),\qquad
 h^\bullet(\mathcal T)=(1,1,e_k,e_k).
\]
The program \texttt{e\_push\_exact.py} independently checks these
counts for all $1965$ coprime pairs with $k,l\le80$; the monomial
argument proves them for all pairs.

On the normalization $\nu:E=\mathrm{dP}_6\to W\subset X$,
the branchwise normal line is $K_E$: the invariant cusp form
$d\log x_1\wedge d\log x_2\wedge dq$ trivializes $K_X$
on the entire cusp neighbourhood.  Local normal coefficients
therefore give the exact sequence
\[
 0\longrightarrow\mathcal T\longrightarrow T_X
 \longrightarrow\nu_*K_E\longrightarrow0,
 \qquad h^\bullet(K_E)=(0,0,1).
\]
We must determine its connecting map when $k>1$.
The dual logarithmic sequence, tensored with $K_X$, has horizontal
subline $\bar f^*\mathcal O(13z-2d)$ and vertical quotient
$\bar f^*(\mathcal O(z-d)\oplus\mathcal O(2z-d))$.
The latter has no global sections, by the same equations
$3A+4B=1,2$ above.  Consequently
\[
 H^0(\Omega_X^1(\log W)\otimes K_X)
 =H^0(\mathfrak B,\mathcal O(13z-2d)).
\]
For $k>1$ this is generated by the monomial
$x^{3l-2}y^{k-2}$ multiplying the horizontal logarithmic
differential.  Both $x$ and $y$ are nonzero over $p_0$, so its
residue on $E$ is nonzero.  The ordinary logarithmic residue
sequence, tensored with $K_X$, identifies this residue map with
the Serre dual of
$H^2(E,K_E)\to H^3(X,\mathcal T)$.
One can check this identification directly in each chart
$q=z_0z_1z_2$: the residues of $dq/q$ are one on each branch,
and contraction with the invariant three-form is the normal
coefficient pairing.  Thus that connecting map has rank $e_k$.
The long exact sequence proves $h^2(T_X)=1$ and $h^3(T_X)=0$;
degrees zero and one are unchanged.

Finally $\partial_c\Pi=e\otimes\gamma$ gives the constant
groupoid cocycle $\lambda\mapsto\gamma(\lambda)e$ on every
regular torus.  On the enlarged finite roots the derivative of the
affine generator is $(r_j/m_j)e$, and the derivative of its
punctured connector is $(\log s_j)/(2\pi i)\,r_je$.
The identity $m_j(r_j/m_j)=\gamma(w_j)$ verifies the power
relation, while $\gamma A_j=\gamma$ verifies conjugation with
every lattice arrow.  Changing $\log s_j$ by $2\pi i N$
adds the derivative of the integral arrow $Nw_j$; the connector
formula is used only on the punctured overlap.  Its central
root-chart representative is the constant $(r_j/m_j)e$,
which is holomorphic.  At the cusp $\gamma$ kills the toric periods,
so the constant deck cocycle is exactly its extending cocycle
above.  Its image under the two isomorphisms
$H^1(\mathcal V)\simeq H^1(\mathcal T)\simeq H^1(T_X)$
is $e\otimes[\gamma]\ne0$.  The chain rule gives the asserted
$u$-class.  The Kuranishi local algebra is
$\mathbb C\{t\}/I$ since $h^1=1$.  The existing family induces
$t\mapsto g(c)$ with $g'(c_0)\ne0$.  The analytic inverse
function theorem makes this substitution injective; every
element of $I$ maps to zero, so $I=0$.  Its classifying map is
then an isomorphism of germs, proving local miniversality.
\end{proof}

\begin{proposition}[E.2.2 (Proved): local separation and canonical dimension]
The multiplicity pair $(3k,4l)$ is constant on every sufficiently
small deformation of one of the preceding members.  Distinct
pairs therefore determine disjoint open loci among these
isomorphism classes.  Moreover
\[
 \kappa(X)=\begin{cases}-\infty,&kl<3k+4l,\\1,&kl>3k+4l.
 \end{cases}
\]
Equality $kl=3k+4l$ never occurs under $\gcd(3k,4l)=1$.
\end{proposition}
\begin{proof}
Local miniversality identifies every small deformation with the
fixed-data period family; E.1.2 makes its multiplicity pair an
intrinsic biholomorphic invariant.  The canonical formula of
E.2.1 and connected fibres identify its pluricanonical sections
with those of $\mathcal O(r(kl-3k-4l))$ on the root curve.
Negative weight has no sections.  Positive weight has linearly
growing section counts along multiples of $12kl$, giving
Kodaira dimension one.  Equality would imply $k\mid4l$;
coprimality forces $k=1$, and then $l=3+4l$ is impossible.
\end{proof}

\begin{proposition}[E.2.3 (Proved): separation in irreducible analytic families]
Suppose a proper holomorphic submersion over a connected
irreducible reduced complex base has two fibres of the type
in E.2.1.  Then their multiplicity pairs agree, even if some
other fibres have not been constructed by this method.
For these members the asymptotic number
\[
 s(X)=\lim_{r\to\infty}
 \frac{h^0(X,K_X^r)-h^0(X,K_X^{-r})}{r}
 =\frac1{12}-\frac1{3k}-\frac1{4l}
\]
recovers the ordered multiplicity pair uniquely.
\end{proposition}
\begin{proof}
For coprime $m,n$, the number of nonnegative solutions of
$mi+nj=N$ is $N/(mn)+O(1)$ as $N\to\infty$:
there is one permitted residue class of $j$ modulo $m$, and
its representatives lie between $0$ and $N/n$.
The canonical formula in E.2.1 proves the displayed limit,
with a negative value when the anticanonical counts grow.
Write $1/12-s(X)=a/b$ in lowest positive terms.
Since $\gcd(m+n,mn)=1$, we have $a=m+n$, $b=mn$.
Thus $m,n$ are the two roots of $T^2-aT+b$.
Their order is determined because only $m$ is divisible by
three and only $n$ is divisible by four.

For every fixed integer $r$, apply the finite-free-complex
theorem of \cite[Chapter III, Theorem 4.1]{BS76push} to
the relative canonical power, which is an invertible flat
sheaf for a holomorphic submersion.  Its zeroth cohomology
dimension has analytic jumping loci.  On an irreducible
base any value constant on a nonempty open set is its
generic value: a proper analytic jumping locus has empty
interior.  Local miniversality in E.2.1 makes these dimensions
constant near either of the two given fibres, since all
nearby fibres have the same multiplicities and their
canonical formula is independent of the period parameter.
Consequently the two fibres have equal dimensions for
every positive and negative $r$.  Their limits $s(X)$
are equal, and the preceding arithmetic recovers the same
pair.  This proof uses relative cohomology only; it makes
no general invariance assertion for Kodaira dimension.
\end{proof}

\begin{problem}[E.2.4 (Open): arbitrary connected deformation chains]
We have not excluded a connected reducible deformation base,
or a chain of irreducible families, joining different pairs
through smooth members outside this period construction.
E.2.3 rules out a single irreducible analytic family containing
two such pairs.  At an unconstructed intervening fibre,
plurigenera may jump on two different branches, so the proof
does not turn $s(X)$ into a deformation invariant on the
entire space of complex structures.  E.2.2 also does not
make these open loci closed.  The attempted local-miniversality
argument fails at that closedness step.  E.2.5 below proves
that a one-sided smooth specialization has a holomorphic
central algebraic reduction, and E.2.6 removes the limiting
fixed divisor and degree loss when the canonical slope is
positive.  A horizontal fixed divisor for negative slope,
and loss of the fractional canonical sections on a different
deformation branch through an intervening fibre, remain
outside those persistence statements.
The outer boundary obstruction of Part D does not control
the parameter end at zero.  A proof of full separation
would require such persistence at every unconstructed
smooth boundary, or an invariant that remains constant
along an arbitrary connecting deformation chain.
\end{problem}

% Owner: boundary_d.  Separate from the E-owned master source.
% The assertions below concern one-sided smooth specializations, not
% arbitrary deformation chains.  The latter remain E.2.4.
\subsubsection*{Smooth specializations of the canonical pencil}

\begin{proposition}[E.2.5 (Proved): a fractional canonical pencil at a smooth limit]
Let $\pi:\mathcal X\to\Delta$ be a proper holomorphic submersion
whose fibres are diffeomorphic to $S^6$.  Suppose that for every
$t\ne0$ the fibre is one of the members of E.2.1, with the same
coprime multiplicities $(3k,4l)$.  Set
\[
 A=12kl,\qquad D=kl-3k-4l\ne0.
\]
There is a unique holomorphic line bundle $\mathcal L$ on
$\mathcal X$, up to isomorphism, satisfying
\[
 \mathcal L^{D}=K_{\mathcal X/\Delta}^{A}.
\]
Its restriction to $X_t$ for $t\ne0$ is the coarse pencil
$f_t^*\mathcal O_{\mathbb P^1}(1)$.  The coherent sheaf
$\pi_*\mathcal L$ is locally free of rank two, and its central
fibre injects into $H^0(X_0,\mathcal L_0)$.

The central manifold has algebraic dimension one and a
holomorphic algebraic reduction $g_0:X_0\to\mathbb P^1$ with
connected fibres.  If $s_0,s_1$ are the two central sections
obtained from a local frame of $\pi_*\mathcal L$, there are a
nonconstant rational map $R:\mathbb P^1\to\mathbb P^1$ of
degree $e\ge1$ and an effective Cartier divisor $E_0$ such that
\[
 [s_0:s_1]=R\circ g_0,\qquad
 \mathcal L_0\simeq g_0^*\mathcal O(e)\otimes\mathcal O(E_0).
\]
Here $[s_0:s_1]$ denotes the holomorphic extension of the ratio;
the displayed divisor is its common fixed divisor.  The
statement does not assert that $E_0=0$.
\end{proposition}
\begin{proof}
By Ehresmann's theorem the smooth bundle over the contractible
disc is differentiably trivial.  Consequently
$H^1(\mathcal X,\mathbb Z)=H^2(\mathcal X,\mathbb Z)=0$.
The exponential sequence identifies
\[
 \operatorname{Pic}(\mathcal X)=H^1(\mathcal X,\mathcal O_{\mathcal X})
\]
as groups.  This is a complex vector space, so multiplication by
the nonzero integer $D$ is an automorphism.  This proves the
existence and uniqueness of the \emph{relative} root, including
when $D<0$.  The same argument on each fibre makes its Picard
group uniquely divisible.  E.2.1 gives
$K_{X_t}^{A}=(f_t^*\mathcal O(1))^{D}$, so uniqueness identifies
$\mathcal L_t$ with this actual pencil, not just with a numerical
class.

Proper direct-image coherence gives $\mathcal V=\pi_*\mathcal L$.
It is torsion-free: a local section killed by a nonzero function
of $t$ vanishes on the inverse image of an open set, hence
vanishes as a section of the invertible sheaf on the smooth
total space.  A coherent torsion-free sheaf on a smooth curve
is locally free.  Its rank is two because the noncentral
algebraic reductions have connected fibres.  Applying $\pi_*$
to
\[
 0\longrightarrow\mathcal L\xrightarrow{\,t\,}\mathcal L
 \longrightarrow i_*\mathcal L_0\longrightarrow0
\]
shows directly that $\mathcal V/t\mathcal V$ injects into
$H^0(X_0,\mathcal L_0)$.  Thus the two central sections are
linearly independent even if central cohomology jumps; no
surjectivity of base change has been assumed.

We use precisely the non-holomorphic-meromorphic-map theorem
of \citet[Theorem~2.1]{CDP2020D}, also written out in
\citet[Theorem~3.1 and Corollary~3.2]{LRS2018D}.
For a compact complex threefold with $b_2=0$, a nonconstant
meromorphic map to a curve which fails to be holomorphic
would force $c_3=0$.  Since $c_3(X_0)[X_0]=2$, every
meromorphic function on $X_0$ extends as a holomorphic map
to $\mathbb P^1$.  This argument uses neither the different
holomorphic-fibration assertion of CDP, Theorem~2.2, nor its
Corollary~2.3.  Those assertions cannot be applied here.

In particular the ratio of $s_0,s_1$ gives a nonconstant
holomorphic map $h_0:X_0\to\mathbb P^1$.  Also $a(X_0)\le1$:
otherwise two algebraically independent meromorphic functions
would give a holomorphic map $X_0\to\mathbb P^1\times\mathbb P^1$
with two-dimensional image.  At a point of differential rank
two, after choosing finite affine target coordinates $u,v$
centred at its image, the global meromorphic function $u/v$
would have the local indeterminacy of the quotient of two
independent coordinates.  It cannot extend holomorphically
there.  This contradicts the preceding theorem.  Hence
$a(X_0)=1$.

Take the Stein factorization $h_0=R\circ g_0$.  Its intermediate
space is a smooth compact curve.  Connectedness of the fibres
injects its $H^1(\ ,\mathbb Z)$ into $H^1(X_0,\mathbb Z)=0$
by the degree-one Leray sequence, so it is $\mathbb P^1$.
Every meromorphic function on $X_0$ factors rationally through
$g_0$: the image of its holomorphic map together with $g_0$
is an analytic curve in $\mathbb P^1\times\mathbb P^1$;
normalizing this curve gives a finite factor of $g_0$, whose
degree is one because $g_0$ has connected fibres.  Thus $g_0$
is the algebraic reduction, and $R$ has some degree $e\ge1$.

Finally, on a target chart where the first homogeneous
coordinate of $h_0$ is nonzero, $s_1=h_0s_0$ in a local
trivialization, and the common ideal is generated by $s_0$.
The analogous assertion holds on the other chart.  These local
generators define a nonzero section of
$\mathcal L_0\otimes h_0^*\mathcal O(-1)$ and an effective
Cartier zero divisor $E_0$.  This gives the stated factorization
of $\mathcal L_0$ with no residual codimension-two base locus.
\end{proof}

\begin{proposition}[E.2.6 (Proved): positive canonical slope preserves the fibration]
Under the hypotheses of E.2.5, suppose $D>0$.  Then $E_0=0$
and $e=1$.  After shrinking the disc and choosing a frame of
$\pi_*\mathcal L$, its evaluation defines a proper flat
holomorphic map
\[
 g:\mathcal X\longrightarrow\mathbb P^1\times\Delta
\]
which is the intrinsic algebraic reduction on every fibre.
Its smooth general central fibre is a complex two-torus.

In particular, two families of the type in E.2.5 with positive
canonical slope cannot have the same smooth central manifold
unless their multiplicity pairs agree.  A positive-slope and
a negative-slope family cannot have a common smooth central
manifold either.  These conclusions concern a common smooth
specialization; they do not establish E.2.4 for arbitrary
connected chains.
\end{proposition}
\begin{proof}
Choose a very general smooth fibre $F$ of $g_0$, avoiding the
images of all vertical components of $E_0$ and the jumping
sets for the pluricanonical direct images below.  Its normal
bundle in $X_0$ is trivial.  Adjunction gives
$K_F=K_{X_0}|_F$, and $c_1(K_F)=0$ integrally because
$H^2(X_0,\mathbb Z)=0$.  In particular $F$ is minimal: a
$(-1)$-curve would have canonical intersection $-1$.
Restriction of the root identity in E.2.5 gives
\[
 K_F^{A}\simeq\mathcal O_F(D E_0|_F).
\]
The right hand side is effective since $D>0$, so
$\kappa(F)\ge0$.

We next show $\kappa(F)\le0$, without assuming that the fibres
of an algebraic reduction have algebraic dimension zero.
If $\kappa(F)>0$, some $r>0$ has $h^0(F,K_F^r)\ge2$.
Our very-general choice ensures that this is the generic rank
of the coherent sheaf $g_{0*}K_{X_0}^r$.  A sufficiently
positive twist of this sheaf on $\mathbb P^1$ is globally
generated.  Two of its global sections can therefore be
chosen to be linearly independent on a general fibre.  By
projection formula they are sections of
$K_{X_0}^r\otimes g_0^*\mathcal O(N)$; their ratio is a
meromorphic function on $X_0$ which is nonconstant on $F$.
This contradicts the algebraic-reduction assertion of E.2.5.
The choice of a very general fibre is legitimate: for each
$r$ the exceptional set for the generic cohomology rank is
a proper analytic subset of the base curve, and there are
only countably many $r$.

Thus $\kappa(F)=0$.  The classification of minimal compact
complex surfaces makes $K_F$ torsion; the relevant consequence
is also explicitly recorded in the proof of
\citet[Proposition~3.1.2]{AbashevaDeev2024AB}.
An effective divisor whose line bundle is torsion is zero:
raising its canonical section to a torsion-killing power
gives a holomorphic function on a compact connected manifold,
and that function is a nonzero constant.  Therefore
$E_0|_F=0$.  Every horizontal component of $E_0$ would meet
such a general fibre in a nonempty effective divisor.
Consequently $E_0$ has only vertical components, supported
over finitely many points of $\mathbb P^1$.

We now use connectedness of actual divisor limits to remove
both this vertical fixed part and the possible degree $e$.
Choose a linear combination $s=as_0+bs_1$ for which its
moving central divisor is the union of the $e$ distinct
smooth fibres over $R^{-1}([-b:a])$, all disjoint from the
support of $E_0$.  Also choose it to be a regular fibre on
one specified noncentral member.  Both requirements omit
only finitely many choices in the projective line of
combinations.  Extend it by the corresponding combination
of the frame of $\pi_*\mathcal L$.

Its total zero divisor $\mathcal D\subset\mathcal X$ is a
relative effective Cartier divisor: the section is not
identically zero on any fibre, by independence of the two
sections, including at zero.  In local coordinates, $t$ is
a non-zero-divisor modulo its defining equation.  Hence
$\mathcal D\to\Delta$ is proper and flat.  Its general
noncentral fibre is a smooth connected torus.  There are no
vertical irreducible components of $\mathcal D$, since such
a component would have to be the whole irreducible
threefold $X_t$.  It follows that $\mathcal D$ is reduced
and irreducible: its unique generic component has multiplicity
one, and a Cartier divisor in a smooth manifold has no
embedded components.  Normalize $\mathcal D$ and take the
Stein factorization over $\Delta$.  The finite factor has
degree one over the punctured generic locus; normality of
the disc makes it the disc itself.  Its central fibre, and
therefore the image $\mathcal D_0$, is connected.

But the support of
\[
 \mathcal D_0=E_0+g_0^*R^*([-b:a])
\]
is a disjoint union of the $e$ chosen smooth connected fibres
and, if nonempty, the support of $E_0$ over other base points.
Connectedness forces $E_0=0$ and $e=1$.

The central two sections now generate $\mathcal L_0$ at
every point.  The noncentral two sections already generate
their pencils, so the evaluation
$\pi^*\pi_*\mathcal L\to\mathcal L$ is surjective after
shrinking.  It defines the asserted map $g$.  It is proper
because $\pi$ is proper.  On each $X_t$ the nonconstant
map to the smooth curve $\mathbb P^1$ is flat: its local
structure sheaf is torsion-free over the local discrete
valuation ring of the curve.  The total source is flat
over $t$, so the local criterion for flatness applied to
the parameter $t-t_0$ gives flatness over $(z,t)$.
Equivalently, the local fibre equations form a regular
sequence in the smooth source.  At a regular
point of the central base curve, properness makes $g$ a
proper submersion on a neighbourhood of the corresponding
compact fibre.  Ehresmann's theorem identifies that central
surface differentiably with nearby complex two-tori.
The classification of compact complex surfaces, or its
special case for the differentiable four-torus, then makes
the central surface a complex two-torus as well.

Suppose now two positive-slope families have the same
central manifold $X_0$.  Its algebraic reduction is intrinsic
up to $\operatorname{PGL}_2$, so the preceding result makes
both central pencil lines the same non-torsion line
$L=g_0^*\mathcal O(1)$.  With indices $i=1,2$ the identities
$K_{X_0}^{A_i}=L^{D_i}$ give
$L^{D_1A_2-D_2A_1}=\mathcal O$.  The line $L$ is non-torsion
because a positive power has a nonzero section with zeros.
Thus $D_1/A_1=D_2/A_2$, and the exact recovery of the ordered
pair from this slope in E.2.3 gives the assertion.

Finally, a common smooth limit from opposite signs would
have a nonzero section of a positive power of $K_{X_0}$
and of a negative power, by upper semicontinuity (or by
the same injective direct-image argument).  Raising them
to a common power, their nonzero product is a constant;
both are nowhere zero.  Hence $K_{X_0}$ is torsion.
The exponential sequence and $H^1(X_0,\mathbb Z)=
H^2(X_0,\mathbb Z)=0$ make $\operatorname{Pic}(X_0)$
torsion-free, so $K_{X_0}$ would be trivial.  Its nonzero
holomorphic three-form $\Omega$ is closed and exact because
$H^3(X_0,\mathbb C)=0$.  Stokes' theorem would give
$\int_{X_0}\Omega\wedge\overline\Omega=0$, whereas its
appropriate constant multiple is a positive volume form.
This contradiction excludes the common limit.
\end{proof}

\begin{remark}[Scope after E.2.6]
E.2.5 removes the possibility of a genuinely non-holomorphic
central meromorphic reduction.  E.2.6 removes its fixed
divisor and degree loss when $D>0$.  For $D<0$ the restriction
identity instead permits an effective antipluricanonical
divisor on a Hopf-type general surface; the preceding proof
does not exclude that possibility.  Even in the positive
case, a different local deformation branch through an
unconstructed limit might lose the sections of the
fractional canonical line.  Neither local closure of all
the known loci nor separation along arbitrary reducible
bases or arbitrary deformation chains has been proved.
Thus E.2.4 remains open, with these two precise restrictions
on a possible connecting chain.
\end{remark}


\subsection{E.3. Analytic foundations and the historical first version}
\begin{proposition}[E.3.1 (Known, with chartwise derivation): base change and duality]
For the proper flat root-curve morphisms in E.2.1, the canonical
geometric-point comparison
\[
 (R^q\bar f_*\mathcal O_X)\otimes k(b)
 \longrightarrow H^q(X_b,\mathcal O_{X_b})
\]
is an isomorphism for $q=0,1,2$, and the direct images are locally
free.  The trace induces the canonical isomorphism
\[
 (R^2\bar f_*\mathcal O_X)^\vee
 \simeq\bar f_*\omega_{X/\mathfrak B}.
\]
Both comparisons descend on the analytic root-curve groupoid.
These conclusions, and hence the use made of them in R4, require
no result specific to the historical first version of
\cite{DGKXv1push}.
\end{proposition}
\begin{proof}
Take the atlas consisting of the smooth base complement and the
two complete root discs.  The pullback over the $j$th root disc
is the actual torus family before its free affine quotient:
\[
 [J'_j/\mu_{m_j}]
 \mathbin{\times_{[\Delta_{s_j}/\mu_{m_j}]}}
 \Delta_{s_j}=J'_j.
\]
The atlas map, unlike $s_j\mapsto s_j^{m_j}$ on coarse spaces,
is etale.  Properness, smoothness at these roots, and the
flat equation $q=z_0z_1z_2$ at the remaining cusp were checked
in E.2.1.  The atlas source is smooth; all its fibres are local
complete intersections of pure dimension two.  It is therefore
a proper flat Gorenstein family of complex spaces over a reduced
smooth curve.  Torus cohomology and the explicit cusp resolution
give the constant dimensions $(1,2,1)$ on every chart, including
its centre.

Apply B\u anic\u a and St\u an\u a\c sil\u a
\cite[Chapter III, Theorem 4.12(ii), pp.~134--135]{BS76push}:
properness, base flatness of the coherent sheaf, a reduced base,
and constant fibre cohomology give cohomological flatness and
local freeness.  Their Theorem 4.10, pp.~132--133, supplies the
canonical analytic base-change comparison.  On the double atlas
overlap both pullbacks are identified with the direct image of
the same pulled-back family.  These comparisons commute with a
third pullback by their construction, so they satisfy the
groupoid cocycle condition.  The resulting descended vector
bundle is $R^q\bar f_*\mathcal O_X$, and its geometric fibre
comparison is the displayed one.  No assertion about the
cohomology of the reduced multiple surface is substituted for
the torus geometric fibre of the stack map.

Let $p:Y\to U$ denote any of these atlas families.  Proper analytic
duality is the canonical trace identity in
Ramis--Ruget--Verdier \cite[introduction, equation (1),
Section 1, pp.~264--265, and Proposition 2.1, p.~266]{RRV71push}.
The cited article does not number its main identity as a theorem;
these are its actual locations.  The smooth-space dualizing
complex is $K_Y[3]$, and the relative Gorenstein dualizing
complex is $\omega_{Y/U}[2]$; the relevant absolute statements
are Ramis--Ruget \cite[Section 3, Theorem 1, p.~81;
Proposition 2 and the complete-intersection calculation, p.~91]{RR70push}.
Cancelling $p^*K_U[1]$ in proper duality gives
\[
 Rp_*\omega_{Y/U}[2]
 \simeq R\mathcal Hom_U(Rp_*\mathcal O_Y,\mathcal O_U).
\]
The direct images just obtained are locally free, so the positive
sheaf Ext terms in the hyper-Ext spectral sequence vanish.
Degree $-2$ is therefore
$p_*\omega_{Y/U}\simeq(R^2p_*\mathcal O_Y)^\vee$.
The trace is local and invariant under the local biholomorphisms
on the atlas overlaps.  The two pullbacks of this isomorphism
thus agree with the trace for the same pulled-back family and
satisfy the same cocycle condition.  This proves descent.

For the original multiplicities $(3,4)$, E.2.1 specializes to
$\mathcal E=\mathcal O\oplus\mathcal O(-1)$,
$\omega_{X/\mathfrak B}=\bar f^*\mathcal O(1)$ and
$R^2\bar f_*\mathcal O_X=\mathcal O(-1)$.
The nondegenerate cusp cup product identifies
$\det R^1\bar f_*\mathcal O_X$ with this last line, and the
nonzero $[\gamma]$-subline and $H^1(\mathcal O(1))=0$ give
$R^1=\mathcal O\oplus\mathcal O(-1)$.
This derives exactly the direct images used in R4.
\end{proof}

Grauert's original direct-image theorem
\cite[Section 6, Hauptsatz I, and Section 7, Satz 5]{Grauert60push}
was also opened together with its correction \cite{Grauert63push}.
The correction explicitly warns that the exceptional locus in
Section 7, Satz 2 need not be closed.  E.3.1 uses the
constant-cohomology theorem on every chart and the later
base-change formulation, so it does not make that erroneous
closedness inference.  For a generic argument over a smooth
curve one can instead use the finite free complex of
\cite[Chapter III, Theorem 4.1]{BS76push}; its rank-jump minors
give the required local analytic exceptional sets.

\begin{proposition}[E.3.2 (Proved as a source audit): the five historical claims]
For the standard twist vectors, each of the five named theorems
in the attached version 1 of \cite{DGKXv1push} has a complete
argument relative to the explicitly constructed family and the
standard analytic theorems it cites.  We located no failed step
in those five arguments.  Their assertions are not present in
the attached version 2 \cite{DGKXv2push}.  This textual
difference carries no inference about their mathematical truth
or about the reason for their omission.
\end{proposition}
\begin{proof}
We give the dependency audit, rather than infer correctness from
the occurrence of the theorem labels.

For \texttt{p2:thm:tangent-cohomology}, the proof first constructs
the root lift as an actual Cartesian quotient, proves properness
and flatness on its atlas, and computes the cusp normalization
resolution with the multiplicative cup pairing.  The two lemmas
\texttt{p2:lem:analytic-grauert} and
\texttt{p2:lem:analytic-relative-duality} explicitly cite the
standard results checked in E.3.1 and perform descent; they do
not assume an unproved general stack duality principle.  Its
weighted cohomology and the normal sequence then give
$(1,1,1,0)$.  Its algebraic-to-analytic Picard comparison invokes
the applicable proper-stack GAGA theorem
\cite[Theorem 7.4]{PY16push}; the direct analytic Picard
argument in E.2.1 provides another justification.  Our general
calculation reduces to precisely its four acyclic weights
$-1,-2,-5,-6$ at $k=l=1$.  The extra logarithmic cohomology
for $k>1$ is outside that source's assertion.

For \texttt{p2:thm:kuranishi-aut}, the generating-arrow formulas
are $\gamma(\lambda)e$ on a lattice arrow, $0$ on a core deck
arrow, $e/3,-e/4$ on the two affine root generators, and
$\gamma(\lambda)e$ on a cusp arrow.  Its proof verifies the
relations, uses logarithms only on punctured connectors, and
uses the extending toric logarithmic vector on the cusp.
The Leray edge isomorphism carries this class to
$e\otimes[\gamma]\ne0$.  Completeness of the Kuranishi family
then gives a substitution $t\mapsto g(s)$ with $g'(0)\ne0$;
injectivity forces the entire defining ideal, including any
nilpotents, to vanish.  This proves a smooth reduced germ and
miniversality, without incorrectly claiming uniqueness or
universality.  The same theorem proves only
$\operatorname{Aut}^0=\mathbb C^*$: its flow has kernel
$\mathbb Z$, verified on the normalization of the cusp.  Full
automorphisms are treated separately below.

For \texttt{p2:thm:intrinsic-fibration}, adjunction and the
local finite characters give $K^{-2}=f^*\mathcal O_B(1)$.
Connected fibres and the projection formula identify all its
sections with the two base sections.  The complete pencil
therefore forces a biholomorphism to cover a base automorphism.
The three different multiplicity and monodromy types distinguish
the marked values, so this base automorphism is the identity.
This is the same argument as A.1 of the original return; it
needs no assertion that an arbitrary biholomorphism was
initially fibre preserving.

For \texttt{p2:thm:exact-period}, the simultaneous centralizer
is computed entrywise, and the period columns force integral
parameter difference.  Normalization after the genuine cyclic
base change tests the actual affine deck action, giving the
order-four parity obstruction.  Sufficiency is explicit:
choose $F^3=t$, $G^4=t-1$ and
$\ell=(2\pi i)^{-1}\log(F/G)$ on the universal cover, modulo
integral constants.  The source's translation is
\[
 b_{a,h}=-\frac h{12}\Pi_a\widehat\gamma+h\ell e.
\]
Its defects at $g_1,g_2,g_0$ are respectively
\[
 \tfrac h2(\widehat u-\widehat w-\widehat\delta),\qquad
 \tfrac h2(-\widehat w+\widehat\delta),\qquad0.
\]
For even $h$ these are integral.  Removing the known logarithmic
connector on each root chart leaves a holomorphic translation;
at the cusp $\ell=-s/12+$ a holomorphic function cancels the
period's logarithm.  Thus both the map and its inverse extend
over all three fillings.  Appendix calculations give the raw
composition defect $-hk\,e/12$ and the correction
$-h^2e/24$, which makes the action strict.  These identities
were re-executed exactly in
\texttt{code/e\_push\_source\_audit.py}.  The return's torsor
argument is another sufficiency proof, not a repair of a
missing proof in this source.

For \texttt{p2:thm:full-automorphisms}, a self-map's linear
part is first forced to be the identity.  The integral middle
translation cohomology is $\mathbb Z$, with representative
$c_N(g_1)=N(0,1,-2,0)^t$, $c_N(g_2)=0$.
The cusp defect is $N(0,-1,1,0)^t$, whose residual class in
$\Lambda/(\mathbb Z\widehat w+\mathbb Z\widehat\delta)$ is
$(0,-N)$.  A transverse disc through a smooth central point
lifts to the free cusp cover, forcing that residual monodromy
to vanish.  Hence $N=0$.  For the remaining zero class the
local quotient lift gives meromorphic character ratios; their
already single-valued additive logarithms have zero winding,
so the ratios are units and their logarithms extend.  Finite
averaging gives extension on both complete root discs.  The
result is a global section of the vertical bundle, hence a
constant multiple of $e$.  Its period kernel is $\mathbb Z$.
This proves $\operatorname{Aut}=\mathbb C^*$, including the
absence of disconnected components.  It corresponds to the
integral-cocycle and cusp-extension steps of B.4.1 in the return.

The audit program records attachment SHA-256 hashes and all five
label locations in version 1; every label and the corresponding
Kuranishi, Kodaira--Spencer, parity-classification and tangent
statements are absent from version 2.  We separately opened the
version-specific arXiv records and treat these as preprint
versions, not as a journal publication verified by its title.
\end{proof}

\begin{proposition}[E.3.3 (Proved as a source audit): limits of the overlap]
The five audited version-1 results concern the standard
$(3,4)$ twist data.  That text does not classify general twist
vectors, compute the exact maximum $M$, construct a filling as
the family parameter tends to zero, compute a class in
$\Theta_7$, or settle the boundary questions in Parts A--D.
\end{proposition}
\begin{proof}
Its finite generators use only
$\widehat\gamma+2\widehat u-4\widehat w$ and
$-\widehat\gamma-3\widehat u+3\widehat w$ throughout the
deformation and classification arguments.  The cusp completion
constructed there fills the elliptic-base puncture for a fixed
period parameter; it is not a filling at zero of the parameter
$w$.  Its period domain is an admissible region, without the
maximization or the interval certificate of R3.  The concluding
orbit groupoid, quadratic correction, and the Oka results do not
supply the missing parameter degenerations or their topology.
Thus none of these omitted conclusions can be read into the
five theorems merely from the words ``family'' or ``cusp''.
\end{proof}

\subsection{E.4. An independent interval certificate}
\label{sec:e4}

We write $\eta_D$ for the Dedekind eta function and set
\[
 h=z-(1-\tau)/2,\qquad
 \mathcal N^2(z,\tau)=
 \left|\frac{\vartheta_3(h,\tau)}{\eta_D(\tau)}\right|^2
 \exp\left(-\frac{2\pi(\Im h)^2}{\Im\tau}\right).
\]
The half-period theta identity identifies this with the norm used in
R3. The coefficient of $\Im c$ in $D_c$ is one. Consequently the
identity from the supplied return which we use is
\[
 D_{\rm part}=\frac6\pi
 \log\left(\frac{\mathcal N}{2\sqrt3}\right).
\]
It follows directly by taking the imaginary part of the theta formula
for $\beta^{\rm part}$; its Gaussian factor removes
$6(\Im\mu)^2/\Im\tau$. Thus no numerical normalization is inferred
from the new computation. The analytic formula and the first
certificate are recorded in the supplied return, Statements C.1.1
and C.2.3 \citep{family2026}. The result below supplies independent
arithmetic evidence for that same conclusion, without a priority claim.

\begin{theorem}[E.4.1 (Proved): independent certification of every box]
For all 19701 rational rectangles in the supplied certificate, and for
each $\tau$ in the rectangle, both zeros of $\wp_\tau$ satisfy
$\mathcal N^2<2$. The new largest certified upper bound is less than
\[
 1.999877571767410525.
\]
The 19701 rectangles, 36 local rectangles and 174 excluded rectangles
have disjoint interiors and cover exactly
$[0,1/2]\times[17/20,2]$. Every local rectangle is contained in
$|\tau-i|<3/20000$, and every excluded rectangle satisfies
$|\tau|<1$ throughout.
\end{theorem}

\begin{proof}
The programs \path{code/e4_arb_certificate.py} and
\path{code/e4_cover_and_local.py} use respectively 192-bit and 256-bit
Arb arithmetic through python-flint 0.9.0. The arithmetic implementation
is independent of mpmath. We import neither mpmath nor any routine
from the previous certificate. The only numerical data taken from that
certificate are its rational rectangle endpoints. We recompute roots,
preconditioners, radii, contraction bounds, function values and tails.
Arb's midpoint-radius representation and its outward error bounds are
described by \citet[Sections 2--3]{Johansson2016}; the precise
constructor and endpoint semantics used here are documented in
\citet{PythonFlint2026}. The final comparison is made with an upper
endpoint, before conversion to a decimal string.

For a rectangle $Q$, let $z_0$ and $A\ne0$ be the exact dyadic complex
numbers selected by the midpoint Newton iteration. That iteration is
only a choice of constants. For $Z=z_0+[-r,r]+i[-r,r]$ we independently
enclose $\wp_\tau(z_0)$ and $\wp'_\tau(Z)$ for $\tau\in Q$.
We accept only if
\[
 \begin{split}
 &\left|-A\wp_\tau(z_0)
 +(1-A\wp'_\tau(Z))(Z-z_0)\right|_\infty<r,\\
 &|\Re(1-A\wp'_\tau(Z))|
       +|\Im(1-A\wp'_\tau(Z))|<1.
 \end{split}
\]
Here absolute values on enclosures mean their outward upper bounds.
The derivative enclosure and integration along a line segment show
that $z\mapsto z-A\wp_\tau(z)$ is a strict contraction of $Z$ into
itself, uniformly for $\tau\in Q$. Thus it has a genuine zero in $Z$.
The two zeros of $\wp$ are negatives modulo the period lattice, counted
with multiplicity. The theta transformation laws show that
$\mathcal N$ is invariant under these operations; hence one root
enclosure controls both zeros. No continuity or success of Newton
iteration is used in that conclusion.

We use the normally convergent cosine expansion
\[
 \wp_\tau(z)=\frac{\pi^2}{\sin^2\pi z}-\frac{\pi^2}{3}
 +8\pi^2\sum_{n\geq1}\frac{nq^n}{1-q^n}
                      (1-\cos 2\pi nz),\qquad q=e^{2\pi i\tau},
\]
and differentiate it in $z$ for $\wp'$. This identity follows by twice
differentiating the normally convergent Jacobi product and inserting
$E_2=1-24\sum_{n\geq1}nq^n/(1-q^n)$; this also verifies the
Fourier formula used in the supplied return. The new program sums
32 terms directly in cosine form. Put $r_q\geq|q|$,
$H_z\geq|\Im z|$, and $R_q=r_q e^{2\pi H_z}<1$.
For $j=1,2$, the inequality $(N+k)^j\leq(N+1)^jk^j$ gives
\[
 \begin{aligned}
 U_1(a)&=\frac{(N+1)a^{N+1}}{(1-a)^2},\\
 U_2(a)&=\frac{(N+1)^2a^{N+1}(1+a)}{(1-a)^3},
 \qquad
 \sum_{n>N}n^j a^n\leq U_j(a).
 \end{aligned}
\]
We therefore add complex rectangles enclosing discs of radii
\[
 \frac{8\pi^2}{1-r_q}\bigl(U_1(r_q)+U_1(R_q)\bigr),
 \qquad \frac{16\pi^3}{1-r_q}U_2(R_q)
\]
to the two partial sums. These estimates follow from
$|\cos z|,|\sin z|\leq e^{|\Im z|}$ and
$|1-q^n|\geq1-r_q$. They differ from the moment formulas used in
the old implementation.

For $\vartheta_3$ we sum
$1+2\sum_{n=1}^{12}e^{\pi i n^2\tau}\cos(2\pi nh)$.
With $T_*\leq\Im\tau$ and $H_*\geq|\Im h|$, the remaining
absolute value is at most
\[
 \frac{2e^{-169\pi T_*+26\pi H_*}}
 {1-e^{-27\pi T_*+2\pi H_*}}.
\]
The ratio in this expression is tested to be less than one. For eta
we independently sum its logarithm,
\[
 \log\eta_D(\tau)=\frac{\pi i\tau}{12}
          +\sum_{n=1}^{24}\log(1-q^n)+\mathcal E,
 \qquad
 |\mathcal E|\leq
 \frac{r_q^{25}}{(1-r_q)(1-r_q^{25})}.
\]
The logarithms are holomorphic on $|q^n|<1$. The tail follows by
summing $|\log(1-z)|\leq|z|/(1-|z|)$.
We enclose $\mathcal N^2$ as
\[
 ((\Re\vartheta_3)^2+(\Im\vartheta_3)^2)
 \exp\left(-2\Re\log\eta_D-
                   \frac{2\pi(\Im h)^2}{\Im\tau}\right),
\]
and test its upper endpoint against the exact integer 2.

All 19701 tests succeeded. The full output records each new enclosure;
the largest is enclosed by
\[
 1.999877571767410524070536478244687423578792301873142913
       \ \pm 9.26\cdot10^{-56}.
\]
This is the enclosure of a certified upper bound, not an estimate of
the actual maximum. Its box has input index 17866, with zero-based
indexing. The file digest of the input is recorded in the output.

The cover check is exact rational arithmetic. At all 438 distinct
vertical edges we maintain the active closed vertical intervals. In
every open vertical slab they form an ordered partition of
$[17/20,2]$, each with multiplicity one. The last active set is empty;
all leaves lie in the root rectangle, and the sum of their exact areas
is $23/40$. The closed rectangles also cover their common edges.
The omitted and excluded conditions are checked separately using
rational endpoint comparisons. Finally,
$2(1/10000)^2<(3/20000)^2$ puts the entire omitted square inside the
analytic disc used below. This verifies both the inequalities and
the coverage, rather than just replaying the old list of successful
flags.
\end{proof}

\begin{theorem}[E.4.2 (Proved): quantitative local maximum]
If $|\tau-i|\leq3/20000$ and $\wp_\tau(z)=0$, then
$\mathcal N(z,\tau)\leq\sqrt2$, with equality precisely at
$\tau=i$. The proof uses the explicit constants below.
\end{theorem}

\begin{proof}
Set $x=\tau-i$, $G(h,\tau)=\wp_\tau((1-\tau)/2+h)$ and
$H(h,\tau)=\log(\vartheta_3(h,\tau)/\eta_D(\tau))$.
Both functions are even in $h$. On $|h|\leq1/5$, $|x|\leq1/10$,
the point $(1-\tau)/2+h$ has distance at least $0.9/2-0.2$ in
imaginary part from any lattice point, so $G$ has no pole.
Put
\[
 A_k=\sum_{n\geq1}n^ke^{-0.9\pi n^2}\cosh(0.4\pi n),
 \quad a_0=2A_0,\ a_1=4\pi A_1,\ a_2=8\pi^2A_2,
 \quad r=e^{-1.8\pi}.
\]
Then $a_0<0.225$, so the theta series has no zero there and a
holomorphic logarithm exists. From
$G=-\partial_h^2\log\vartheta_3-\pi^2E_2/3$ we obtain
\[
 \begin{aligned}
 |G|&\leq\frac{a_2}{1-a_0}+
 \frac{a_1^2}{(1-a_0)^2}+
 \frac{\pi^2}{3}\left(1+\frac{24r}{(1-r)^3}\right)<18.373<20,\\
 |H|&\leq-\log(1-a_0)+\frac{1.1\pi}{12}
                  +\frac{r}{(1-r)^2}<0.547<0.6.
 \end{aligned}
\]
The new certified upper bounds are respectively
$18.372504032305989$ and $0.546229724979279$.
For the sums $A_k$ we retain ten terms. If $n=11+j$, then
$n^2\geq121+23j$ and $n^k\leq11^k(1+j)^k$; replacing cosh by
the positive exponential bounds the tail by a geometric series with
ratio $e^{-20.3\pi}$ and its first two moment sums. These formulas,
including all outward rounding, are in the accompanying program.

At $(0,i)$ the coefficients are
\[
 G=0,\qquad G_{h^2}=-g_2(i)/4,\qquad G_x=id,
 \qquad H=\log\sqrt2,\quad H_x=0,\quad H_{h^2}=-\pi/2,
\]
where $d>0$. The first two identities follow from the half-period
value of $\wp$ and $\wp''=6\wp^2-g_2/2$. If
$e_2(\tau)=\pi^2(\vartheta_2^4-\vartheta_4^4)/3$, differentiating
its Gaussian series gives $e_2'(i)=id$. We compute
\[
 \begin{split}
 10.03&<d<10.04,\\
 47.26&<g_2(i)/4<47.30.
 \end{split}
\]
More precisely the certified enclosures begin
$d=10.03061933746602249879\ldots$ and
$g_2(i)/4=47.26818003230846307326\ldots$.
The computation of $d$ uses 16 terms in the two Gaussian series and
their derivatives, followed by geometric Gaussian tail bounds.
For $g_2/4$ it uses
\[
 \frac{g_2(i)}4=\frac{\pi^4}{3}
 \left(1+240\sum_{n\geq1}\frac{n^3e^{-2\pi n}}{1-e^{-2\pi n}}\right)
\]
with 24 terms and the bound
$\sum_{n>24}n^3q^n/(1-q^n)
\leq25^3q^{25}(1+4q+q^2)/(1-q)^5$.
The modular transformation of $\vartheta_3/\eta_D$ gives $H_x=0$;
that of eta gives $E_2(i)=3/\pi$. Inserting this in the displayed
formula for $G$ gives $H_{h^2}=-\pi/2$.

Take $\epsilon=3/20000$, $r_h=3/100$,
$s=\epsilon/(1/10)=3/2000$ and $r_2=(r_h/(1/5))^2=9/400$.
Cauchy's coefficient estimates give on $|h|=r_h$
\[
 |G(h,i)|\geq
 \left(47.26-\frac{500r_2}{1-r_2}\right)r_h^2
       >0.03217594373401534,
\]
while
\[
 |G(h,i+x)-G(h,i)|\leq
 \frac{20s}{(1-s)(1-r_2)}<0.03073664204747019.
\]
Rouch\'e's theorem therefore encloses both zeros in $|h|<r_h$.
The comparison with the total number of zeros of an elliptic function
of order two identifies these with all roots modulo the lattice.

Write $G=e_2(\tau)+h^2Q(h,\tau)$. The same Cauchy bounds imply
\[
 \left|Q+\frac{g_2(i)}4\right|
 \leq\frac{500(r_2+s/(1-s))}{1-r_2}<12.277368,
 \quad
 |e_2(\tau)-idx|\leq\frac{2000|x|^2}{1-s}.
\]
At a root, first bounding the left side of the last inequality below
and then bounding it above yields respectively
\[
 |x|<6.123343|h|^2<7|h|^2,
 \qquad |h|^2<0.000044338219<\frac1{20000}.
\]
For example the two rational majorants actually evaluated are
\[
 \frac{47.30+Q_*}{10.03-2000\epsilon/(1-s)},
 \qquad
 \frac{10.04\epsilon+2000\epsilon^2/(1-s)}{47.26-Q_*},
\]
where $Q_*=500(r_2+s/(1-s))/(1-r_2)$; both denominators are
strictly positive.

In the even power series of $H$, put
$\alpha=|x|/(1/10)\leq70|h|^2$ and
$\beta=|h|^2/(1/5)^2=25|h|^2$. After removal of its displayed
constant and linear coefficients, the remaining bivariate series is
bounded by
$0.6(\alpha^2+\alpha\beta+\beta^2)/((1-\alpha)(1-\beta))$.
The change from $1/\Im\tau$ to one in the Gaussian contributes
at most $7\pi|h|^4/(1-\epsilon)$. Thus a bound for the total
coefficient of $|h|^4$ is
\[
 \frac{0.6(70^2+25\cdot70+25^2)}
 {(1-s)(1-25/20000)}+\frac{7\pi}{1-\epsilon}
 <4399.024<5000.
\]
Since $\Re(-\pi h^2/2)-\pi(\Im h)^2=-\pi|h|^2/2$, we have
\[
 \log\mathcal N-\log\sqrt2
 \leq-\frac\pi2|h|^2+5000|h|^4
 <-\left(\frac\pi2-\frac14\right)|h|^2
\]
whenever $h\ne0$. The coefficient is greater than $1.32079632679$.
If $h=0$, the lower bound for $e_2$ forces $x=0$. This proves the
claim with all constants controlled by a second arithmetic engine.
\end{proof}

\begin{corollary}[E.4.3 (Proved): confirmation of the maximum in R3]
For the normalization in the prompt, $M=-(3/\pi)\log6$, attained
only at $t=1$. In particular the radius 216 has an independent
arithmetic certificate.
\end{corollary}
\begin{proof}
The theta transformation formulas preserve $\mathcal N$ under modular
transport, lattice translation, sign, and
$(z,\tau)\mapsto(\bar z,-\bar\tau)$; hence the half modular domain
suffices. Theorems E.4.1--E.4.2 cover its portion $\Im\tau\leq2$.
For $T=\Im\tau\geq2$, choose $|\Im z|\leq T/2$ by lattice
translation. With $b=43/100$ and $r=e^{-4\pi}$, if
$|\Im z|\geq b$ the Fourier formula gives
\[
 \left|\frac{\wp_\tau(z)}{\pi^2}+\frac13\right|
 \leq\frac1{\sinh^2(\pi b)}+
 8\left\{\frac r{(1-r)^3}
 +\frac1{2(1-r)}\left(
 \frac{e^{-2\pi}}{(1-e^{-2\pi})^2}+
 \frac{e^{-6\pi}}{(1-e^{-6\pi})^2}\right)\right\}<\frac13.
\]
Thus a zero has $|\Im z|<b$. The function
$\cosh(\pi y)e^{-\pi y^2/T}$ is increasing for $0\leq y\leq b$:
concavity and the positive endpoint value of $\tanh(\pi y)-y$
give its nonnegative logarithmic derivative. The Jacobi product now
implies
\[
 \mathcal N\leq2\cosh(\pi b)
 e^{-\pi/3-\pi b^2/2}
 \exp\left(\frac{2\cosh(2\pi b)r}{1-r}
                +\frac{r^2}{1-r^2}\right)<1.082<\sqrt2.
\]
This bound decreases with $T\geq2$. Both strict inequalities are
independently checked in \path{code/e4_cover_and_local.py}.
At $\tau=i$, $\vartheta_3(0,i)/\eta_D(i)=\sqrt2$, so equality
does occur. The period identity at the beginning of the subsection
then gives the claimed value of $M$ and
$e^{-\pi M}=6^3=216$. The invariant $j(\tau)=1728t$ sends its
unique modular equality orbit to $t=1$.
\end{proof}

\subsection{E.5. The two-parameter family and its actual singular fibres}
\begin{proposition}[E.5.1 (Refuted as stated; corrected): the input configuration]
The configuration in the question is not the configuration of
\cite{ViaTwo26push}.  The elliptic surface
$y^2=x^3+tx+1$ has fibre $III^*$ at infinity.  The constructed
threefold has two rank-one fibres at $p_1,p_2$, a rank-two
hexagonal fibre at $p_3$, and the multiple bielliptic fibre
$4B_*$ at infinity.
\end{proposition}
\begin{proof}
Put $s=t^{-1}$, $X=s^2x$, and $Y=s^3y$.  The minimal equation is
$Y^2=X^3+s^3X+s^6$, with coefficient and discriminant orders
$(3,6,9)$; this is the $III^*$ equation, as in
\cite[Lemma 2.1]{ViaTwo26push}.  The discriminant at finite
points is $-16(4t^3+27)$ with three simple zeros.  In the
threefold's multiplicative periods the orders at these points
are respectively $(1,0),(1,0),(1,1)$:
\[
 h(z,v)=(q\alpha(q)z,\beta(q)v),\qquad
 k(z,v)=(\gamma(q)z,\lambda q^\epsilon\delta(q)v),
\]
where all four coefficients are units, $\epsilon=0$ at $p_1,p_2$,
and $\epsilon=1$ at $p_3$.  Thus the first two torus ranks
drop by one and the last by two.  The explicit matrices of
E.5.2 independently have nilpotent ranks $1,1,2$; their exact
calculation rules out the assertion that all three are rank one.
These are the types stated in Proposition 3.1 and Theorem 1.1
of the cited source, not a correction to that source.
\end{proof}

\begin{proposition}[E.5.2 (Proved): monodromy, centralizer, and integral lattices]
Use the source's left-action marking
$H^1(F,\mathbb Z)=\mathbb Z\langle\psi,u,w,\delta\rangle$
and the ordered meridians $*,1,2,3$ with product one.
The cohomological matrices are
\[
\begin{aligned}
 T_*&=\begin{pmatrix}1&0&-2&-1\\0&0&-1&0\\0&1&0&-1\\0&0&0&1\end{pmatrix},&
 T_1&=\begin{pmatrix}1&0&2&2\\0&1&1&1\\0&0&1&0\\0&0&0&1\end{pmatrix},\\
 T_2&=\begin{pmatrix}1&0&0&0\\0&1&0&0\\0&-1&1&0\\0&0&0&1\end{pmatrix},&
 T_3&=\begin{pmatrix}1&0&0&-1\\0&1&1&0\\0&0&1&0\\0&0&0&1\end{pmatrix}.
\end{aligned}
\]
On ordinary dual homology put $H_i=(T_i^{-1})^t$,
$\ell=(1,-1,-1,0)^t$ and $\widehat\delta=(0,0,0,1)^t$.
Then
\[
 \operatorname{Cent}_{M_4(\mathbb Q)}(H_*,H_1,H_2,H_3)
 =\mathbb Q I\oplus\mathbb Q E_{41},
\]
and the cohomological centralizer is
$\mathbb Q I\oplus\mathbb Q E_{14}$.
The fixed and norm lattices at the multiple fibre are
\[
 \begin{aligned}
 \Lambda^{H_*}&=\mathbb Z\ell\oplus\mathbb Z\widehat\delta,\\
 N_4\Lambda&=4\mathbb Z\ell\oplus2\mathbb Z\widehat\delta,\\
 \Lambda^{H_*}/N_4\Lambda&\simeq\mathbb Z/4\oplus\mathbb Z/2.
 \end{aligned}
\]
The primitive vanishing lattices at the three reduced fibres are
\[
 K_1=\mathbb Z(\widehat w+\widehat\delta),\qquad
 K_2=\mathbb Z\widehat u,\qquad
 K_3=\mathbb Z\widehat w\oplus\mathbb Z\widehat\delta.
\]
Their invariant lattices are
\[
 \begin{split}
 \Lambda^{H_1}&=\mathbb Z(\widehat\psi-2\widehat u)
       \oplus\mathbb Z\widehat w\oplus\mathbb Z\widehat\delta,\\
 \Lambda^{H_2}&=\mathbb Z\widehat\psi\oplus
       \mathbb Z\widehat u\oplus\mathbb Z\widehat\delta,\\
 \Lambda^{H_3}&=K_3.
 \end{split}
\]
\end{proposition}
\begin{proof}
The matrices are Proposition 4.2 of \cite{ViaTwo26push}, whose
derivation uses the elliptic monodromies, the winding of $P$,
and the primitive rank-two cusp.  In particular it determines
the initially undetermined winding integer by the saturation
index $d^2=1$, rather than assuming this winding from the
desired topology.  Direct multiplication gives
$T_*T_1T_2T_3=I$, $T_*^4=I$ and
$(T_i-I)^2=0$ for $i=1,2,3$.
For clarity the homological matrices used in our calculation are
\[
\begin{aligned}
 H_*&=\begin{pmatrix}1&0&0&0\\-2&0&-1&0\\0&1&0&0\\1&1&0&1\end{pmatrix},&
 H_1&=\begin{pmatrix}1&0&0&0\\0&1&0&0\\-2&-1&1&0\\-2&-1&0&1\end{pmatrix},\\
 H_2&=\begin{pmatrix}1&0&0&0\\0&1&1&0\\0&0&1&0\\0&0&0&1\end{pmatrix},&
 H_3&=\begin{pmatrix}1&0&0&0\\0&1&0&0\\0&-1&1&0\\1&0&0&1\end{pmatrix}.
\end{aligned}
\]
Solving the sixty-four rational equations $CH_i=H_iC$
gives $C=aI+bE_{41}$ with free $a,b$; the retained exact
row-reduction program displays a basis and verifies its
dimension two.  Inverse transpose gives the stated
cohomological centralizer.  Solving $(H_*-I)v=0$ gives
$v=(a,-a,-a,d)^t$ integrally, and
\[
 N_4=I+H_*+H_*^2+H_*^3
 =\begin{pmatrix}4&0&0&0\\-4&0&0&0\\-4&0&0&0\\0&2&-2&4\end{pmatrix}.
\]
In the basis $(\ell,\widehat\delta)$ its columns are
$(4,0),(0,2),(0,-2),(0,4)$, giving exactly the displayed
norm group.  Each $K_i$ is the actual integral image of
$H_i-I$; its Smith form has only unit nonzero entries.
The kernels of these three matrices give the remaining fixed
lattices.  All products, rational row reductions, Hermite forms,
and Smith forms are executed in
\texttt{code/e\_push\_two\_parameter.py}.
\end{proof}

We now specify the parameter domain, because the source proves
a family on a sufficiently small bidisc, not on the entire
$\mu$-plane.  Set
\[
 S_\mu:\ y^2=x^3+tx+\mu t+1,\qquad
 P_\mu=(-\mu,\sqrt{1-\mu^3}),
\]
with the square root equal to one at zero.  Label the three roots
$p_i(\mu)$ of $4t^3+27(\mu t+1)^2$ holomorphically near zero;
$p_3$ is the zero of the chosen linearization.
Choose a holomorphic generator of its one-dimensional space of
sections vanishing at $p_3$ and scale it by $\lambda$.
For the added homology generator take the inverse of the
contracting lift.  On the universal cover of the four-punctured
base its period columns then have the form
\[
 \Pi_{\mu,c}=
 \begin{pmatrix}2z_\mu&\tau_\mu&1&0\\
 \beta_\mu+c&z_\mu&0&1\end{pmatrix},
 \qquad \lambda=e^{-2\pi i c},\qquad
 \partial_c\Pi=e\otimes\psi.
\]
Here $z_\mu$ is a lifted elliptic logarithm in the chosen marking;
changing the generator of the linearization changes
$\beta_\mu$ and the origin of $c$ by a holomorphic function of
$\mu$, and changes neither the integer-difference criterion
below nor the original $\lambda$-labelled bidisc.
This choice of inverse lift fixes the sign of $c$.
At $p_3$ its second period has jump $-1$, exactly canceled by
$H_3\widehat\psi=\widehat\psi+\widehat\delta$.
The covariance identities for all four matrices, including
$\tau_*=-1/\tau$, $z_*=(z-1)/\tau$,
$\beta_* =\beta-1-2(z-1)^2/\tau$, are checked symbolically
in the same program.

\begin{theorem}[E.5.3 (Proved): classification on the specified parameter domain]
There are $\epsilon>0$ and a real constant $C$ such that the
source construction, with its fixed logarithmic vector $\ell$,
defines a proper holomorphic family
$X_{\mu,c}$ over
$\{ |\mu|<\epsilon\}\times\{\operatorname{Im}c<C\}$.
It has the singular fibres of E.5.1 and the markings above.
Two of its members are biholomorphic if and only if
\[
 \mu'=\mu,\qquad c'-c\in2\mathbb Z.
\]
Every such biholomorphism covers the identity of the elliptic
base.  The quotient parameter in the second variable is
$v=e^{-\pi i c}$, with $\lambda=v^2$.
On a sufficiently small bidisc centred at any
$(0,\lambda_0)$ in Theorem 1.1 of \cite{ViaTwo26push}, choose
one logarithm of $\lambda$.  Within that actual bidisc,
\[
 X_{\mu,\lambda}\simeq X_{\mu',\lambda'}
 \quad\Longleftrightarrow\quad
 (\mu,\lambda)=(\mu',\lambda').
\]
\end{theorem}
\begin{proof}
The equations of Proposition 6.4 in \cite{ViaTwo26push}
give the simultaneous elliptic family, line bundle, linearization,
and good-reduction comparison.  On a smaller closed $\mu$-disc
all unit factors in its four local fillings are uniformly bounded
above and below.  Its contraction estimates therefore have a
common nonzero bound $|\lambda|<\eta$.  Pulling back by
$\lambda=e^{-2\pi i c}$ gives the indicated half-plane.
The integral markings and the four-torsion section $\ell/4$
continue uniquely on this simply connected parameter domain.
On every relatively compact subset the source's finite root
quotients and the fixed rank-one and hexagonal fans give
holomorphic charts with uniform bounds.  Their comparisons are
the specified section-preserving ones, so they agree on
intersections and define a family.  Properness is local on
the parameter domain and follows from the compact four-piece
construction there; smoothness follows from the same free
finite action and unimodular charts.  We make no assertion
about continuation to arbitrary $\mu$.

The intrinsic fibration is the anticanonical pencil.  To verify
the needed input from Proposition 6.1 of the source, on the
smooth core the line-coordinate form $dv/v$ wedged with a
base two-form identifies $K_X$ with $f^*\mathcal O(-1)$.
The forms $dq\wedge d\log z\wedge d\log v$ extend without
zeros or poles in both rank-one charts $q=ab$ and rank-two
charts $q=abc$.  At the finite root at infinity the source's
coordinate change gives
\[
 ds\wedge\frac{dX}{Y}\wedge\frac{dv}{v}
 =4\,dr\wedge\frac{dx'}{y'}\wedge\frac{dv'}{v'},
 \qquad s=r^4.
\]
It is invariant under the actual affine deck map and descends
through its free quotient.  Thus the same canonical formula
holds everywhere.  Connected fibres and projection formula
give $H^0(-K_X)=f^*H^0(\mathcal O(1))$, without base points.
A biholomorphism therefore induces an affine automorphism of
the elliptic base, fixing its distinguished multiple point
at infinity and its distinguished hexagonal point $p_3$.
It may initially interchange the two rank-one points.

Put $Q(\mu)=(p_1-p_3)/(p_2-p_3)$.
The source computes $p_i'(0)=-9/(2p_i(0))$.
If $Q'(0)$ vanished, the three velocities would have the form
$Ap_i(0)+B$; the polynomial $Az^2+Bz+9/2$ would vanish
at three distinct numbers, a contradiction.
Also $Q(0)^2-Q(0)+1=0$, so $Q(0)\ne Q(0)^{-1}$.
Shrink $\epsilon$ so that $Q$ is injective and its image
is disjoint from the image of $Q^{-1}$.  Preservation or
interchange of the rank-one values gives respectively
$Q(\mu')=Q(\mu)$ or $Q(\mu')=Q(\mu)^{-1}$.
The second possibility is excluded, and the first gives
$\mu'=\mu$.  The base automorphism now fixes three distinct
points individually and is the identity.

On regular fibres a biholomorphism is affine, and its integral
linear map commutes with all four $H_i$.  E.5.2 gives
$aI+bE_{41}$.  Comparing the constant third and fourth
period columns makes its complex linear part $aI$;
unimodularity gives $a=\alpha\in\{1,-1\}$.
The first column then gives $h=c'-c\in\mathbb Z$ and
lattice map $\alpha(I-hE_{41})$.
Normalize after the actual fourth-root base change.  Its
affine generator is the one with fourth power $\ell$, not
an arbitrary generator of its cyclic subgroup.  Summing the
translation defect around this generator gives
\[
 (I+hE_{41})\ell-\alpha\ell
 =(1-\alpha)\ell+h\widehat\delta\in N_4\Lambda.
\]
The norm lattice in E.5.2 forces $\alpha=1$ and $h$ even.
This proves necessity, including the multiple fibre's full
non-reduced structure.

We prove sufficiency with all four collar conditions.
Let $h$ be even and put
\[
 b_*=-\tfrac h4(\widehat u+\widehat w),\qquad
 L_*=-\tfrac h2\widehat u.
\]
Direct multiplication gives
\[
 (I-H_*)b_*=\tfrac h4\widehat\delta+L_*.
\]
Consequently $\zeta\mapsto\zeta+\Pi_{\mu,c}b_*$ on the
complete root torus family conjugates the two affine actions,
modulo the integral arrow $L_*$.  It and its inverse extend
over $r=0$ and descend.  On the punctured root collar its
middle translation includes
$h(\log r)/(2\pi i)e$; this is the derivative-free collar
term for the change of the invariant translation
$\Pi\ell/4$, not a logarithm used on the central chart.

At a rank-one point the vanishing coordinate is the first
toric coordinate in its chain cover, while the other
direction is the universal cover of the elliptic factor.
An extending translation has integral winding only in $K_j$.
After removing the corresponding power of $q$, its toric
coefficient is a holomorphic unit and its elliptic translation
extends holomorphically.  These are the full analytic collar
conditions.  Necessity follows by lifting a transverse disc
through a smooth central point to this free cover: its residual
defect in $\Lambda/K_j$ is zero.  The coordinate ratios on
the cover then give the asserted unit and holomorphic
remainder.  Conversely an integral power of $q$ translates
the infinite chain by an integer number of components;
the holomorphic remainders give torus and elliptic translations.
These maps and their inverses extend at both branches of every
node.  At the hexagonal point the same description uses two
toric windings in $K_3$, and integer translation of the fan
extends over its three double curves and both triple points.
This criterion includes regularity of the remainder; membership
of the defect in $K_j$ alone would not exclude an essential
singularity of a preassigned punctured translation.

Choose the three integral defects
\[
 L_1=0,\qquad L_2=\tfrac h2\widehat u,\qquad
 L_3=-\tfrac h2\widehat w.
\]
They belong to $K_1,K_2,K_3$, respectively, and are realized
by the preceding toric translations with initially constant
unit remainder.  In primitive cocharacter coordinates these
are $1$, $q_2^{h/2}$, and $(q_3^{-h/2},1)$, where the
monodromy defect of $q^a$ is defined as $a$ times the chosen
vanishing generator.  The required sphere relation is
\[
 L_*+H_*L_1+H_*H_1L_2+H_*H_1H_2L_3=0.
\]
It holds as an exact integer identity, checked in the program.
Thus the four local choices have compatible additive defects
and define a torsor of allowed middle translations.

We verify the vanishing of its analytic obstruction, including
the difference sheaf.  The root curve has just the order-four
point at infinity and is $\mathfrak C=\mathbb P(1,4)$,
with $K_{\mathfrak C}=\mathcal O(-5)$ and coarse pullback
$\mathcal O_B(1)=\mathcal O_{\mathfrak C}(4)$.
Vertical logarithmic fields on a rank-one chart are generated
by $a\partial_a-b\partial_b$ and the elliptic vector;
on a hexagonal chart they are the two height-zero toric fields.
The root chart has the ordinary vertical tangent frame.
Their linear transition functions depend only on the base;
they define a rank-two bundle $\mathcal E$ on $\mathfrak C$.
The fixed line-coordinate field supplies a subline
$\mathcal O\subset\mathcal E$ everywhere in these logarithmic
frames.  The logarithmic determinant sequence and cancellation
of the three reduced fibre divisors give
$K_X=\bar f^*(K_{\mathfrak C}\otimes(\det\mathcal E)^{-1})$.
The verified canonical formula and
$\bar f_*\mathcal O_X=\mathcal O_{\mathfrak C}$ therefore
give $\det\mathcal E=\mathcal O(-1)$.
The extension splits since
$H^1(\mathfrak C,\mathcal O(1))=0$, by the same two-chart
Laurent calculation with weights $1,4$.
The coarse difference sheaf is $V_B=\rho_*\mathcal E$:
on the root chart this is precisely finite equivariance,
and on each reduced collar it is precisely the holomorphic
remainder just described.  Finite averaging is exact, so
\[
 H^1(B,V_B)=H^1(\mathfrak C,\mathcal E)
 =H^1(\mathcal O\oplus\mathcal O(-1))=0.
\]
The torsor's Cech class vanishes.  Correcting the four local
translations by its holomorphic zero-cochain gives a global
middle translation and exact equality with all four local
extensions.  Their inverses have the opposite allowed
translations and glue as well.  This constructs the required
biholomorphism over $B$.

Finally choose a single logarithm of $\lambda$ on the source's
small bidisc.  Its $c$-values can be placed in a strip of
width less than two.  The criterion just proved then says
$c'=c$ and $\mu'=\mu$, which is equivalent to equality of
the original two parameters.  Conversely equality gives the
identity, with the fixed root, rank-one and hexagonal charts.
\end{proof}

The matrix data and the independent-variation argument are in
\cite[Proposition 4.2 and Proposition 6.4]{ViaTwo26push}; the
source states dimension at least two and does not state the
integer classification on the logarithmic lift.  We derived
the latter here from the marked norm and collar calculations.
For E.1--E.5 our searches on 15 November 2026 included the
exact arXiv identifiers together with ``classification'',
``Kuranishi'', ``deformation'', and ``parameters'', and the
references and version-specific records of the attached
papers.  No further paper resolving the general deformation
equivalence in E.2.4 was located.  This is a record of the
search performed, not a completeness or priority assertion;
the query responses and opened-source metadata are retained
in \texttt{evidence/E\_search*.json}.
